% ============================================================
%  DRAFT TAMBAHAN BAB 3 (bab3.tex)
%  Sisipkan setiap blok di tempat yang ditandai.
%  Semua konten dibuat berdasarkan implementasi nyata di FinalProject/
% ============================================================

% ============================================================
% BLOK A — Sisipkan setelah \subsection{Desain Sistem dan Arsitektur Penelitian}
%           (sebelum \subsection{Pengumpulan dan Simulasi Data})
% ============================================================

\subsection{Perancangan Database dan Penyimpanan Data}

Sistem menyimpan seluruh data nasabah secara persisten menggunakan \textbf{SQLite} dengan nama file nasabah.db yang disimpan di direktori serving/. Database diinisialisasi secara otomatis pada saat \textit{startup} layanan oleh fungsi init\_db() melalui modul database.py, yang mengimpor seluruh baris dari dataset sintetis synthetic\_credit\_risk\_data.csv apabila tabel belum tersedia.

Database terdiri atas dua tabel. Tabel nasabah merupakan tabel utama yang menyimpan 40 lebih atribut setiap nasabah, mulai dari data identitas, profil pekerjaan, kondisi finansial, parameter risiko hasil model, hingga indikator makroekonomi. Tabel sent\_emails berfungsi sebagai log audit yang mencatat setiap pengiriman email konfirmasi kepada pengambil keputusan. Relasi antara kedua tabel ini diikat melalui kolom customer\_id sebagai \textit{foreign key}.

\begin{table}[H]
\centering
\caption{Kelompok Kolom Utama Tabel nasabah}
\label{tab:db-cols}
\begin{tabularx}{\textwidth}{|p{3.5cm}|X|}
\hline
\textbf{Kelompok} & \textbf{Kolom Utama} \\ \hline
Identitas & customer\_id (PK), loan\_id, name\_masked, age, gender, marital\_status \\ \hline
Pekerjaan \& Pendapatan & employment\_type, pension\_status, monthly\_income, salary\_source, income\_stability, salary\_volatility \\ \hline
Kesehatan Finansial & bureau\_score, DTI, past\_due\_months, credit\_history\_length, avg\_account\_balance, other\_active\_loans\_count \\ \hline
Detail Pinjaman & product\_type, loan\_amount, term\_months, interest\_rate, collateral\_flag, ltv, outstanding\_balance \\ \hline
Output Model ML & PD, LGD, EAD, expected\_loss, risk\_grade, missed\_payment\_flag \\ \hline
Status Keputusan & approval\_status (Pending / Accepted / Rejected) \\ \hline
Makroekonomi & macro\_inflation, macro\_unemployment, macro\_interest\_rate, macro\_property\_index \\ \hline
\end{tabularx}
\end{table}

Seluruh operasi \textit{Create, Read, Update, Delete} (CRUD) terhadap tabel nasabah diekspos melalui REST API berbasis \textbf{FastAPI} (app.py). Tabel~\ref{tab:crud-bab3} merangkum pemetaan \textit{endpoint} API ke fungsi database.

\begin{table}[H]
\centering
\caption{Mapping \textit{Endpoint} REST API ke Fungsi Database}
\label{tab:crud-bab3}
\begin{tabularx}{\textwidth}{|p{1.8cm}|p{6cm}|X|}
\hline
\textbf{Metode} & \textbf{Endpoint} & \textbf{Fungsi Database} \\ \hline
GET    & /api/nasabah                          & get\_nasabah\_list() \\ \hline
GET    & /api/nasabah/\{id\}                   & get\_nasabah\_by\_id() \\ \hline
POST   & /api/nasabah                          & create\_nasabah() \\ \hline
PUT    & /api/nasabah/\{id\}                   & update\_nasabah() \\ \hline
DELETE & /api/nasabah/\{id\}                   & delete\_nasabah() \\ \hline
GET    & /api/nasabah/\{id\/approve}           & update\_approval\_status("Accepted") \\ \hline
GET    & /api/nasabah/\{id\/reject}            & update\_approval\_status("Rejected") \\ \hline
POST   & /api/nasabah/\{id\/send-confirmation} & create\_sent\_email() \\ \hline
GET    & /api/sent-emails                      & get\_sent\_emails() \\ \hline
\end{tabularx}
\end{table}

Agar API tahan terhadap ketidakcocokan skema, sebelum setiap operasi INSERT atau UPDATE, modul database.py mengkueri kolom yang tersedia melalui PRAGMA table\_info(nasabah) sehingga hanya kolom yang benar-benar ada di skema yang ditulis. Kolom approval\_status yang tidak terdapat pada dataset CSV asli ditambahkan secara dinamis melalui ALTER TABLE saat init\_db() pertama kali dijalankan.

% ============================================================
% BLOK B — Sisipkan sebagai sub-subsection baru di dalam
%           \subsection{Implementasi MLOps}, setelah sub-bab
%           Deployment Model dan Layanan Inferensi.
%           Letakkan SEBELUM sub-bab Monitoring Drift.
% ============================================================

\subsubsection{Mekanisme Persetujuan Kredit Berbasis Email (\textit{Human-in-the-Loop})}

Sistem tidak memberikan keputusan kredit secara otomatis. Setiap hasil prediksi model harus dikonfirmasi secara manual oleh analis kredit atau pejabat pemutus melalui mekanisme \textit{Human-in-the-Loop} (HITL) berbasis email. Mekanisme ini diimplementasikan melalui alur berikut:

\begin{enumerate}
    \item Setelah model menghasilkan prediksi, analis kredit memicu \textit{endpoint} POST /api/nasabah/\{id\/send-confirmation} dengan menyertakan alamat email reviewer.
    \item Sistem menyusun \textbf{email HTML} yang berisi ringkasan risiko nasabah---mencakup \textit{bureau score}, pendapatan bulanan, jumlah pinjaman, \textit{risk grade}, dan probabilitas \textit{default}---beserta dua tombol interaktif: \textbf{Setuju} dan \textbf{Tolak}.
    \item Email dikirimkan kepada reviewer melalui \textbf{n8n webhook} (/webhook/send-confirmation) yang meneruskan pesan via server SMTP Gmail.
    \item Setiap pengiriman email dicatat secara otomatis ke tabel sent\_emails sebagai \textit{audit trail}.
    \item Klik pada tombol Setuju memanggil GET /api/nasabah/\{id\/approve}, sedangkan klik Tolak memanggil GET /api/nasabah/\{id\/reject}. Kedua \textit{endpoint} ini memperbarui kolom approval\_status di database menjadi "Accepted" atau "Rejected" dan menampilkan halaman konfirmasi kepada reviewer.
    \item Apabila email konfirmasi dikirim ulang untuk nasabah yang sama, sistem secara otomatis mereset approval\_status kembali ke "Pending" agar siklus peninjauan ulang dimulai dari awal.
\end{enumerate}

\begin{figure}[H]
\centering
\begin{tikzpicture}[node distance=1.0cm and 2.0cm]
  \node[brstep] (pd)     {Prediksi PD\\dan risk\_grade selesai};
  \node[brstep, below=of pd] (email)  {Kirim email konfirmasi\\via n8n (SMTP)};
  \node[brdecision, below=of email] (decide) {Keputusan\\Analis?};
  \node[brok, right=3cm of decide] (acc)  {approval\_status\\= \textbf{Accepted}};
  \node[brerr, left=3cm of decide] (rej)  {approval\_status\\= \textbf{Rejected}};
  \node[brstep, above=1.2cm of email] (reset) {Reset ke\\approval\_status\\= \textbf{Pending}};

  \draw[brarrow] (pd)     -- (email);
  \draw[brarrow] (email)  -- (decide);
  \draw[brarrow] (decide) -- node[above,font=\footnotesize]{Setuju} (acc);
  \draw[brarrow] (decide) -- node[above,font=\footnotesize]{Tolak} (rej);
  \draw[brarrow, dashed, bend left=50] (email.east) to
    node[right,font=\footnotesize]{Kirim ulang\\email} (reset.east);
  \draw[brarrow, dashed] (reset) -- (email);
\end{tikzpicture}
\caption{Siklus approval\_status --- keputusan kredit berbasis konfirmasi email manusia.}
\label{fig:hitl-flow}
\end{figure}

Mekanisme ini memastikan bahwa model \textit{machine learning} hanya berperan sebagai \textbf{alat bantu keputusan}, bukan pengganti penilaian manusia. Seluruh keputusan final tetap berada di tangan analis kredit yang berwenang.

% ============================================================
% BLOK C — Ganti atau perluas \subsubsection{Orkestrasi Workflow
%           dan Pelacakan Eksperimen} yang sudah ada.
%           Tambahkan paragraf tentang pipeline kedua n8n.
% ============================================================

% [Letakkan paragraf ini SETELAH paragraf yang sudah ada tentang n8n retraining]

Selain \textit{pipeline} \textit{retraining}, \textit{workflow} n8n juga mengintegrasikan jalur \textbf{ingesti data nasabah berbasis email} secara otomatis sebagai \textit{pipeline} kedua yang terpisah. Setiap jam, n8n memonitor \textit{inbox} Gmail untuk email bersubjek \textit{"Credit Application"} yang menyertakan lampiran berformat PDF. Formulir yang diterima secara otomatis diunduh, diteruskan ke skrip parse\_pdf.py untuk proses ekstraksi data, kemudian disimpan ke database nasabah melalui \textit{endpoint} POST /api/nasabah. Dengan demikian, proses \textit{onboarding} nasabah baru dapat dilakukan sepenuhnya tanpa intervensi manual --- formulir kredit yang dikirimkan melalui email langsung masuk ke database dan siap untuk diprediksi oleh model.

\begin{table}[H]
\centering
\caption{Dua \textit{Pipeline} Utama n8n dalam Sistem MLOps}
\label{tab:n8n-pipelines}
\begin{tabularx}{\textwidth}{|p{3.5cm}|X|p{3.5cm}|}
\hline
\textbf{Pipeline} & \textbf{Alur} & \textbf{Trigger} \\ \hline
\textit{Retraining Pipeline} & Schedule → Jalankan trainingCredit.py → Cek AUC $>$ 0.75 → Promosi ke Production MLflow → \textit{Reload} model & Terjadwal (7 hari sekali) \\ \hline
\textit{Ingestion Pipeline} & Gmail Trigger → Unduh lampiran PDF → parse\_pdf.py → Validasi JSON → POST /api/nasabah & Setiap jam (polling Gmail) \\ \hline
\end{tabularx}
\end{table}

Pelacakan eksperimen dilakukan menggunakan \textbf{MLflow}, yang mencatat parameter \textit{hyperparameter}, metrik evaluasi (AUC, F1-Score), dan artefak model setiap kali proses pelatihan dijalankan. Selain pencatatan eksperimen, sistem juga menggunakan \textbf{MLflow Tracing} (@mlflow.trace) untuk mencatat setiap permintaan inferensi ke \textit{endpoint} /predict beserta versi model yang digunakan. Kemampuan \textit{audit trail} per-prediksi ini mendukung kebutuhan kepatuhan regulasi, termasuk ketentuan transparansi model sesuai POJK Nomor 29 Tahun 2024.

% ============================================================
% BLOK D — Ganti atau perluas \subsubsection{Sistem Grading
%           Risiko Kredit} yang sudah ada.
%           Tambahkan formula bernomor + zona ketiga.
% ============================================================

% [Letakkan setelah paragraf yang menjelaskan grade A-D dan Undefined]

Penerapan batas usia pelunasan dilakukan melalui sistem logika hibrida yang menghitung usia nasabah pada saat jatuh tempo pinjaman. Formulasi yang digunakan adalah sebagai berikut:

\begin{equation}
    \text{Usia Pelunasan} = \text{Usia Saat Pengajuan} + \frac{\text{Tenor (Bulan)}}{12}
    \label{eq:usia-pelunasan}
\end{equation}

Apabila nilai usia pelunasan pada Persamaan~\ref{eq:usia-pelunasan} melebihi 80 tahun, sistem memberikan status Grade Undefined dan menghitung rekomendasi tenor maksimal yang diperbolehkan secara dinamis:

\begin{equation}
    \text{Tenor Maksimal (Bulan)} = (80 - \text{Usia Saat Pengajuan}) \times 12
    \label{eq:tenor-maks}
\end{equation}

Berdasarkan kedua persamaan tersebut, sistem membagi kondisi nasabah ke dalam \textbf{tiga zona keputusan}:

\begin{enumerate}
    \item \textbf{Zona Lolos} --- Usia pelunasan $\leq$ 80 tahun. Data diteruskan ke model XGBoost secara normal untuk memperoleh \textit{risk grade} A, B, C, atau D.
    \item \textbf{Zona Mitigasi Dinamis} --- Usia nasabah saat pengajuan masih di bawah 80 tahun, tetapi tenor yang diminta menyebabkan usia pelunasan melewati 80 tahun. Sistem memberikan Grade Undefined dan menghitung rekomendasi tenor maksimal menggunakan Persamaan~\ref{eq:tenor-maks}. Jika nasabah memiliki asuransi jiwa kredit aktif yang nilainya mencukupi jumlah pinjaman, sistem juga dapat merekomendasikan dispensasi melalui komite kredit.
    \item \textbf{Zona \textit{Joint-Borrower}} --- Usia nasabah saat pengajuan telah mencapai 80 tahun atau lebih. Pada kondisi ini, pengurangan tenor tidak lagi memadai karena sisa batas usia bernilai nol atau negatif. Sistem merekomendasikan alternatif pengajuan berbasis \textit{joint-income} atau \textit{co-borrower} bersama anggota keluarga inti yang berusia di bawah 60 tahun.
\end{enumerate}

Pemisahan tiga zona ini penting karena model \textit{machine learning} hanya mempelajari korelasi statistik antarfitur dan tidak secara eksplisit memahami batas kelayakan usia pelunasan yang diatur oleh kebijakan manajemen risiko institusi. Dengan demikian, logika bisnis pada lapisan API berfungsi sebagai mekanisme pengaman yang melengkapi prediksi model.

% ============================================================
% BLOK E — Ganti atau perluas \subsubsection{Continuous Training
%           dan Retraining Otomatis} yang sudah ada.
%           Tambahkan metrik evaluasi model baru.
% ============================================================

% [Letakkan setelah kalimat pembuka sub-bab yang sudah ada]

% [Letakkan setelah kalimat pembuka sub-bab yang sudah ada]

Keandalan \textit{Continuous Training} diukur bukan dari mendeteksi \textit{drift} pada fitur individu, melainkan langsung pada evaluasi performa model baru hasil \textit{retraining}. Setelah n8n memicu skrip \textit{retraining} (baik karena jadwal rutin, akumulasi data, maupun \textit{trigger} manual), sistem melakukan prosedur \textit{Quality Gate} terhadap model baru menggunakan metrik:

\begin{enumerate}
    \item \textbf{Area Under the Curve (AUC)} --- mengukur kemampuan model secara agregat dalam membedakan kelas \textit{default} dan \textit{non-default}. Model baru harus memiliki nilai AUC minimal 0.75 atau lebih tinggi dari model produksi yang ada saat ini.
    \item \textbf{F1-Score} --- memastikan keseimbangan antara presisi (meminimalkan risiko gagal bayar yang terlewat) dan \textit{recall} (meminimalkan penolakan terhadap nasabah yang sebenarnya berisiko rendah).
\end{enumerate}

Apabila model baru lolos \textit{Quality Gate} ini, skrip secara otomatis mendaftarkan model tersebut ke \textit{MLflow Model Registry} dan memicu \textit{endpoint} khusus pada aplikasi FastAPI untuk melakukan \textit{hot-reload}, sehingga aplikasi langsung melayani prediksi menggunakan model terbaru tanpa mengalami \textit{downtime}.

Agar \textit{Continuous Training} dapat beroperasi secara \textit{closed-loop} pada lingkungan produksi, setiap permintaan prediksi ke \textit{endpoint} /predict secara otomatis dicatat oleh \textit{MLflow Tracing}. Pencatatan ini berisi seluruh fitur input nasabah, probabilitas \textit{default} yang dihasilkan model, dan \textit{timestamp} prediksi. Data operasional historis ini dikumpulkan dan menjadi sumber data \textit{ground truth} ketika label performa kredit aktual nasabah telah diperbarui, yang kemudian digunakan kembali oleh n8n sebagai dataset untuk siklus \textit{Continuous Training} berikutnya. Tanpa mekanisme pencatatan ini, proses \textit{retraining} otomatis tidak akan memiliki suplai data produksi yang relevan.
